\documentclass[11pt]{article}

\usepackage[a4paper,margin=1in]{geometry}
\usepackage{amsmath,amssymb,amsthm,mathtools}
\usepackage[T1]{fontenc}
\usepackage{lmodern}
\usepackage{microtype}
\usepackage{hyperref}

\title{}
\author{}
\date{}

\begin{document}
\maketitle

For \(0<a<1\) and \(n\ge2\), consider the matrix
\[
A_n(a)
=
\operatorname{tridiag}(a,0,1)
=
\begin{bmatrix}
0 & 1 \\
a & 0 & \ddots \\
& \ddots & \ddots & \ddots \\
&& \ddots & 0 & 1 \\
&&& a & 0
\end{bmatrix}\in \mathbb{R}^{n\times n}.
\]
With fixed \(a\), define the real-axis least singular value
\[
\mu_n(x;a)\coloneqq s_{\min}(xI_n-A_n(a)),
\qquad x\in\mathbb{R}.
\]
Let
\[
\lambda_{n,k}(a)
=
2\sqrt a\cos\frac{(n+1-k)\pi}{n+1},
\qquad k=1,\dots,n,
\]
be the eigenvalues of \(A_n(a)\) in increasing order.
Define the \(k\)-th gap height by
\[
\beta_{n,k}(a)
\coloneqq
\max_{\lambda_{n,k}(a)\le x\le\lambda_{n,k+1}(a)}
\mu_n(x;a),
\qquad k=1,\dots,n-1,
\]
and the largest real-axis barrier by
\[
\tau_n(a)\coloneqq\max_{1\le k\le n-1}\beta_{n,k}(a).
\]
Then, for every \(0<a<1\) and every \(n\ge2\),
\[
\tau_{n+1}(a)\le\tau_n(a).
\]

\paragraph{A structural reduction of the problem.}
Let \(R_n\) denote the reversal matrix,
\[
(R_n)_{ij}=\delta_{i,n+1-j}.
\]
Since
\[
A_n(a)^T=R_nA_n(a)R_n,
\]
the matrix
\[
H_n(x):=R_n\bigl(xI_n-A_n(a)\bigr)
\]
is real symmetric. Moreover,
\[
H_n(x)^2
=
R_n\bigl(xI_n-A_n(a)\bigr)R_n\bigl(xI_n-A_n(a)\bigr)
=
\bigl(xI_n-A_n(a)\bigr)^T\bigl(xI_n-A_n(a)\bigr).
\]
Consequently
\[
\boxed{\;
\mu_n(x;a)
=
\min_{\eta\in\sigma(H_n(x))}|\eta|
\;}
\tag{1}
\]
for every real \(x\). In entries,
\[
(H_n(x))_{ij}
=
\begin{cases}
x, & i+j=n+1,\\
-a, & i+j=n,\\
-1, & i+j=n+2,\\
0, & \text{otherwise}.
\end{cases}
\tag{2}
\]
Thus the singular-value problem is exactly the problem of locating the eigenvalue of the real symmetric anti-tridiagonal pencil \(H_n(x)\) closest to zero.

Let
\[
p_0(x)=1,\qquad p_1(x)=x,\qquad
p_{m+1}(x)=xp_m(x)-ap_{m-1}(x)\quad(m\ge1).
\]
Then
\[
p_m(x)=\det(xI_m-A_m(a))
=
a^{m/2}U_m\!\left(\frac{x}{2\sqrt a}\right),
\tag{3}
\]
where \(U_m\) is the Chebyshev polynomial of the second kind. Hence the zeros of \(p_n\) are precisely
\[
2\sqrt a\cos\frac{j\pi}{n+1},\qquad j=1,\dots,n.
\]

For \(x\notin\sigma(A_n(a))\), the inverse of \(xI_n-A_n(a)\) is given explicitly by
\[
\bigl(xI_n-A_n(a)\bigr)^{-1}_{ij}
=
\begin{cases}
\dfrac{p_{i-1}(x)p_{n-j}(x)}{p_n(x)}, & i\le j,\\[8pt]
a^{\,i-j}\dfrac{p_{j-1}(x)p_{n-i}(x)}{p_n(x)}, & i>j.
\end{cases}
\tag{4}
\]
This follows directly by substituting the displayed formula into
\((xI_n-A_n(a))G=I_n\): away from \(i=j\) the entries satisfy the same
three-term recurrence as the \(p_m\), while the jump at \(i=j\) is exactly
the Wronskian identity
\[
p_jp_{n-j}-ap_{j-1}p_{n-j-1}=p_n.
\]

Multiplying (4) on the right by \(R_n\) gives a symmetric matrix
\[
K_n(x):=\bigl(xI_n-A_n(a)\bigr)^{-1}R_n=H_n(x)^{-1}.
\]
Its entries are
\[
(K_n(x))_{ij}
=
\begin{cases}
\dfrac{p_{i-1}(x)p_{j-1}(x)}{p_n(x)}, & i+j\le n+1,\\[10pt]
a^{\,i+j-n-1}\dfrac{p_{n-i}(x)p_{n-j}(x)}{p_n(x)}, & i+j\ge n+1.
\end{cases}
\tag{5}
\]
Because \(H_n(x)\) is symmetric,
\[
\mu_n(x;a)=\frac{1}{\|K_n(x)\|_2},
\qquad x\notin\sigma(A_n(a)).
\tag{6}
\]
Therefore
\[
\beta_{n,k}(a)^{-1}
=
\min_{\lambda_{n,k}(a)<x<\lambda_{n,k+1}(a)}
\|K_n(x)\|_2,
\tag{7}
\]
and
\[
\tau_n(a)^{-1}
=
\min_{1\le k\le n-1}
\min_{\lambda_{n,k}(a)<x<\lambda_{n,k+1}(a)}
\|K_n(x)\|_2.
\tag{8}
\]

There is also an exact congruence diagonalization. Put
\[
\theta_j=\frac{j\pi}{n+1},\qquad
\lambda_j=2\sqrt a\cos\theta_j,
\]
and define the right eigenvectors
\[
v_j(m)=a^{(m-1)/2}\sin(m\theta_j),
\qquad m=1,\dots,n.
\]
Then
\[
A_n(a)v_j=\lambda_jv_j.
\]
Similarly,
\[
w_j(m)=a^{-(m-1)/2}\sin(m\theta_j)
\]
is a left eigenvector. A direct computation gives
\[
R_nv_j=(-1)^{j+1}a^{(n-1)/2}w_j.
\]
Hence
\[
v_\ell^{T}R_nv_j
=
(-1)^{j+1}a^{(n-1)/2}
\sum_{m=1}^{n}\sin(m\theta_\ell)\sin(m\theta_j)
=
(-1)^{j+1}\frac{n+1}{2}a^{(n-1)/2}\delta_{\ell j}.
\]
Writing \(V=[v_1,\dots,v_n]\), one obtains
\[
\boxed{\;
V^TH_n(x)V
=
\frac{n+1}{2}a^{(n-1)/2}
\operatorname{diag}\!\left(
(-1)^{j+1}(x-\lambda_j)
\right)_{j=1}^n.
\;}
\tag{9}
\]
Thus \(H_n(x)\) is a self-adjoint congruence model for the non-normal
Toeplitz resolvent problem.

The monotonicity assertion
\[
\tau_{n+1}(a)\le \tau_n(a)
\tag{10}
\]
would follow from the following localized finite-\(n\) comparison principle:
for every \(n\ge2\), every \(1\le k\le n\), and every \(s>0\),
\[
\{x\in(\lambda_{n+1,k},\lambda_{n+1,k+1}):\mu_{n+1}(x;a)>s\}\ne\varnothing
\]
should imply
\[
\{x\in(\lambda_{n,k-1},\lambda_{n,k}):\mu_n(x;a)>s\}\ne\varnothing
\quad\text{or}\quad
\{x\in(\lambda_{n,k},\lambda_{n,k+1}):\mu_n(x;a)>s\}\ne\varnothing,
\tag{11}
\]
with the evident omission of intervals outside the spectrum.

Indeed, (11) is equivalent to the pointwise gap comparison
\[
\beta_{n+1,k}(a)
\le
\max\{\beta_{n,k-1}(a),\beta_{n,k}(a)\},
\tag{12}
\]
where \(\beta_{n,0}=\beta_{n,n}=0\). Taking the maximum over \(k\) in
(12) gives immediately
\[
\tau_{n+1}(a)
=
\max_{1\le k\le n}\beta_{n+1,k}(a)
\le
\max_{1\le j\le n-1}\beta_{n,j}(a)
=
\tau_n(a).
\]

Thus the original theorem is reduced, without loss, to the finite-dimensional
comparison principle (11). The identities (1)--(9) are unconditional; the
monotonicity (10) requires an additional proof of (11).

\paragraph{Bottleneck analysis and a sharper recursive reduction.}
Put
\[
M_n(x):=xI_n-A_n(a),
\qquad
T_n(x,s):=M_n(x)^TM_n(x)-s^2I_n .
\]
Then
\[
\mu_n(x;a)>s
\quad\Longleftrightarrow\quad
T_n(x,s)\succ0 .
\tag{1}
\]
Thus the desired monotonicity is equivalent to the following level-set implication:
for every \(s>0\),
\[
\{x\in[\lambda_{n+1,1},\lambda_{n+1,n+1}]:T_{n+1}(x,s)\succ0\}\ne\varnothing
\]
implies
\[
\{x\in[\lambda_{n,1},\lambda_{n,n}]:T_n(x,s)\succ0\}\ne\varnothing .
\tag{2}
\]
Indeed, \(\tau_n(a)>s\) holds exactly when the second set in (2) is nonempty.

The matrix \(T_n(x,s)\) is the real symmetric pentadiagonal matrix
\[
T_n(x,s)=
\begin{bmatrix}
d_1& e & a \\
e& d & e & a\\
a&e&d&e&\ddots\\
&\ddots&\ddots&\ddots&\ddots&a\\
&&a&e&d&e\\
&&&a&e&d_n
\end{bmatrix},
\tag{3}
\]
where
\[
d_1=x^2+a^2-s^2,\qquad
d=x^2+a^2+1-s^2,\qquad
d_n=x^2+1-s^2,\qquad
e=-(1+a)x .
\tag{4}
\]
The only difference between the left semi-infinite leading section and the finite matrix
\(T_n(x,s)\) is the final diagonal decrement by \(a^2\).

Define recursively the pivots \(\rho_j=\rho_j(x,s)\) and the modified first off-diagonal
entries \(\gamma_j=\gamma_j(x,s)\) by
\[
\rho_1=x^2+a^2-s^2,\qquad \gamma_1=-(1+a)x,
\tag{5}
\]
and, for \(j\ge2\),
\[
\rho_j
=
x^2+a^2+1-s^2
-\frac{a^2}{\rho_{j-2}}
-\frac{\gamma_{j-1}^2}{\rho_{j-1}},
\tag{6}
\]
\[
\gamma_j
=
-(1+a)x-\frac{a\gamma_{j-1}}{\rho_{j-1}},
\tag{7}
\]
with the convention that the term \(a^2/\rho_{j-2}\) is omitted when \(j=2\).
These are exactly the scalar \(LDL^T\) pivots of the semi-infinite left leading
pentadiagonal section.

For the finite matrix \(T_n(x,s)\), the first \(n-1\) pivots are
\[
\rho_1,\rho_2,\dots,\rho_{n-1},
\]
while the final pivot is
\[
\rho_n-a^2 .
\tag{8}
\]
Therefore Sylvester's criterion gives the exact finite-section test
\[
\boxed{\;
\mu_n(x;a)>s
\quad\Longleftrightarrow\quad
\rho_1(x,s)>0,\dots,\rho_{n-1}(x,s)>0,
\quad
\rho_n(x,s)>a^2 .
\;}
\tag{9}
\]

Equivalently, if \(L_j(x,s)\) denotes the determinant of the \(j\times j\) left leading
semi-infinite pentadiagonal section, with \(L_0=1\), then
\[
\rho_j=\frac{L_j}{L_{j-1}},
\tag{10}
\]
and the finite determinant satisfies the exact identity
\[
\det T_n(x,s)
=
L_n(x,s)-a^2L_{n-1}(x,s).
\tag{11}
\]
Thus
\[
\mu_n(x;a)>s
\quad\Longleftrightarrow\quad
L_1,\dots,L_{n-1}>0,
\qquad
L_n-a^2L_{n-1}>0.
\tag{12}
\]

This is the intrinsic recursive core of the problem.  The monotonicity
\[
\tau_{n+1}(a)\le \tau_n(a)
\tag{13}
\]
is therefore reduced to the following purely scalar continued-fraction assertion:

\[
\boxed{
\begin{minipage}{0.91\linewidth}
For every \(s>0\), if there exists
\(x\in[\lambda_{n+1,1},\lambda_{n+1,n+1}]\) such that
\[
\rho_1(x,s)>0,\dots,\rho_n(x,s)>0,
\qquad
\rho_{n+1}(x,s)>a^2,
\]
then there exists
\(y\in[\lambda_{n,1},\lambda_{n,n}]\) such that
\[
\rho_1(y,s)>0,\dots,\rho_{n-1}(y,s)>0,
\qquad
\rho_n(y,s)>a^2.
\]
\end{minipage}}
\tag{14}
\]
By (9), (14) is exactly the level-set implication (2), hence exactly the desired
monotonicity.

The important point is that (14) is not a routine consequence of principal-submatrix
interlacing.  The obstruction is visible directly from the pivot test.  Positivity of
\(T_{n+1}(x,s)\) gives
\[
\rho_1,\dots,\rho_n>0,\qquad \rho_{n+1}>a^2,
\tag{15}
\]
whereas positivity of \(T_n(x,s)\) would require
\[
\rho_1,\dots,\rho_{n-1}>0,\qquad \rho_n>a^2.
\tag{16}
\]
Thus the missing step is precisely the passage from the final condition
\(\rho_{n+1}>a^2\) to the previous final condition \(\rho_n>a^2\), possibly at a different
real point in the adjacent spectral interval.  Pointwise monotonicity is false; one cannot
hope to prove
\[
\mu_{n+1}(x;a)\le \mu_n(x;a)
\]
for all relevant \(x\).  The proof must instead be a global one-dimensional oscillation
argument for the Riccati recurrence (5)--(7).

The symmetric-pencil formulation gives the same obstruction from another angle.
Let \(R_n\) be the reversal matrix and define
\[
H_n(x):=R_nM_n(x).
\]
Then \(H_n(x)\) is real symmetric and
\[
H_n(x)^2=M_n(x)^TM_n(x),
\]
so
\[
\mu_n(x;a)=\min_{\eta\in\sigma(H_n(x))}|\eta|.
\tag{17}
\]
Moreover \(H_n(x)\) is the affine symmetric pencil
\[
H_n(x)=xR_n-R_nA_n(a).
\tag{18}
\]
The zeroes of the eigenvalue branches of this pencil occur exactly at the eigenvalues of
\(A_n(a)\).  The barrier height in a spectral gap is the largest distance from zero reached
by the closest eigenvalue branch of \(H_n(x)\) between two consecutive zeroes.  At an
interior maximizing point \(x_*\), a corresponding normalized eigenvector \(u\) satisfies
the Hellmann--Feynman stationarity condition
\[
u^TR_nu=0.
\tag{19}
\]
Thus every barrier is produced by a neutral vector of the indefinite metric \(R_n\).
This explains why ordinary Hermitian interlacing after the diagonal similarity
\(A_n(a)\sim \sqrt a\,\operatorname{tridiag}(1,0,1)\) is insufficient: the metric in
which the relevant eigenvalue branches move is indefinite.

The cleanest exact reduction is therefore (14).  A complete proof of the proposed
monotonicity must prove this Riccati comparison.  All preceding identities,
especially the pivot criterion (9), are unconditional and reduce the original matrix
problem to a scalar finite-difference comparison with no loss of information.

\paragraph{Continuation: the real bottleneck in the Riccati comparison.}

I will keep the notation from the preceding reduction.  Thus
\[
M_n(x):=xI_n-A_n(a),\qquad
T_n(x,s):=M_n(x)^TM_n(x)-s^2I_n,
\]
and
\[
\mu_n(x;a)>s
\quad\Longleftrightarrow\quad
T_n(x,s)\succ0.
\tag{1}
\]
The leading-pivot recursion was
\[
\rho_1=x^2+a^2-s^2,\qquad \gamma_1=-(1+a)x,
\tag{2}
\]
and, for \(j\ge2\),
\[
\rho_j
=
x^2+a^2+1-s^2
-\frac{a^2}{\rho_{j-2}}
-\frac{\gamma_{j-1}^2}{\rho_{j-1}},
\tag{3}
\]
\[
\gamma_j
=
-(1+a)x-\frac{a\gamma_{j-1}}{\rho_{j-1}},
\tag{4}
\]
with the term \(a^2/\rho_{j-2}\) omitted for \(j=2\).  The exact finite-section test is
\[
\mu_n(x;a)>s
\quad\Longleftrightarrow\quad
\rho_1(x,s)>0,\dots,\rho_{n-1}(x,s)>0,\qquad
\rho_n(x,s)>a^2.
\tag{5}
\]

The proposed Riccati comparison is therefore the following statement:
\[
\begin{aligned}
&\exists x\in[\lambda_{n+1,1},\lambda_{n+1,n+1}]
\text{ such that }
\rho_1(x,s)>0,\dots,\rho_n(x,s)>0,\quad \rho_{n+1}(x,s)>a^2
\\
&\hspace{1cm}\Longrightarrow
\exists y\in[\lambda_{n,1},\lambda_{n,n}]
\text{ such that }
\rho_1(y,s)>0,\dots,\rho_{n-1}(y,s)>0,\quad \rho_n(y,s)>a^2.
\end{aligned}
\tag{6}
\]

A direct induction on the scalar pivots does \emph{not} prove (6).  The reason is simple and important.  The implication
\[
\rho_{n+1}>a^2
\quad\Longrightarrow\quad
\rho_n>a^2
\tag{7}
\]
is false pointwise.  In fact, in the central gaps one typically has
\[
0<\rho_n(x,s)<a^2
\quad\text{while}\quad
\rho_{n+1}(x,s)>a^2.
\tag{8}
\]
Thus the comparison, if true, must be a global oscillation statement in the real variable \(x\), not a pointwise monotonicity statement in the index \(n\).

The clean way to see the hidden structure is to pass from the scalar pivots to the
self-adjoint singular-value pencil
\[
\mathcal H_n(x):=
\begin{bmatrix}
0&M_n(x)\\
M_n(x)^T&0
\end{bmatrix}.
\tag{9}
\]
Then
\[
\sigma(\mathcal H_n(x))
=
\{-s_1(M_n(x)),\dots,-s_n(M_n(x)),
s_n(M_n(x)),\dots,s_1(M_n(x))\},
\tag{10}
\]
and therefore
\[
\mu_n(x;a)>s
\quad\Longleftrightarrow\quad
\sigma(\mathcal H_n(x))\cap[-s,s]=\varnothing.
\tag{11}
\]
Moreover, \(\mathcal H_n(x)\) is a \(2\times2\) block Jacobi matrix:
\[
\mathcal H_n(x)=
\begin{bmatrix}
B(x)&C\\
C^T&B(x)&C\\
&\ddots&\ddots&\ddots\\
&&C^T&B(x)
\end{bmatrix},
\tag{12}
\]
where
\[
B(x)=
\begin{bmatrix}
0&x\\
x&0
\end{bmatrix},
\qquad
C=
\begin{bmatrix}
0&-1\\
-a&0
\end{bmatrix}.
\tag{13}
\]
Thus the problem is a finite-interval gap problem for a self-adjoint block Jacobi pencil.

The level equations
\[
\pm s\in\sigma(\mathcal H_n(x))
\tag{14}
\]
are equivalent to
\[
\det T_n(x,s)=0.
\tag{15}
\]
The open intervals on which \(T_n(x,s)\succ0\) are precisely the open real gaps between consecutive level-crossing points of the two curves
\[
+s\in\sigma(\mathcal H_n(x)),
\qquad
-s\in\sigma(\mathcal H_n(x)).
\tag{16}
\]
Consequently, the Riccati comparison (6) is exactly the following oscillation statement:

\[
\boxed{
\begin{minipage}{0.91\linewidth}
Whenever the block Jacobi pencil \(\mathcal H_{n+1}(x)\) has a gap around zero of half-width \(s\) at some \(x\) in the spectral hull of \(A_{n+1}(a)\), the preceding block Jacobi pencil \(\mathcal H_n(y)\) has a gap around zero of at least the same half-width at some \(y\) in the spectral hull of \(A_n(a)\).
\end{minipage}}
\tag{17}
\]

This is the nontrivial step.  It is stronger than ordinary Cauchy interlacing.
Indeed, \(\mathcal H_n(x)\) is a principal \(2n\times2n\) block submatrix of
\(\mathcal H_{n+1}(x)\), so Cauchy interlacing gives
\[
\nu_j^{(n+1)}(x)
\le
\nu_j^{(n)}(x)
\le
\nu_{j+2}^{(n+1)}(x),
\qquad j=1,\dots,2n,
\tag{18}
\]
where
\[
\nu_1^{(m)}(x)\le\cdots\le\nu_{2m}^{(m)}(x)
\]
are the ordered eigenvalues of \(\mathcal H_m(x)\).  If
\[
\sigma(\mathcal H_{n+1}(x))\cap[-s,s]=\varnothing,
\tag{19}
\]
then
\[
\nu_{n+1}^{(n+1)}(x)\le -s,
\qquad
\nu_{n+2}^{(n+1)}(x)\ge s.
\tag{20}
\]
But (18) only implies
\[
\nu_n^{(n)}(x)\le \nu_{n+2}^{(n+1)}(x),
\qquad
\nu_{n+1}^{(n)}(x)\ge \nu_{n+1}^{(n+1)}(x),
\tag{21}
\]
which permits both central eigenvalues of \(\mathcal H_n(x)\) to lie inside
\((-s,s)\).  Thus pointwise Hermitian interlacing does not imply (17).

Equivalently, returning to the scalar pivots, positivity of \(T_{n+1}(x,s)\) gives
\[
\rho_1(x,s)>0,\dots,\rho_n(x,s)>0,\qquad \rho_{n+1}(x,s)>a^2,
\tag{22}
\]
whereas positivity of \(T_n(x,s)\) would require
\[
\rho_1(x,s)>0,\dots,\rho_{n-1}(x,s)>0,\qquad \rho_n(x,s)>a^2.
\tag{23}
\]
The final inequality has shifted from \(\rho_{n+1}>a^2\) to \(\rho_n>a^2\), and this shift cannot be justified at the same point \(x\).

There is, however, an exact identity explaining why this obstruction is the only obstruction.  Write
\[
T_{n+1}(x,s)=
\begin{bmatrix}
T_n(x,s)+a^2e_ne_n^T & a e_{n-1}-(1+a)x e_n\\
a e_{n-1}^T-(1+a)x e_n^T & x^2+1-s^2
\end{bmatrix}.
\tag{24}
\]
Hence
\[
T_{n+1}(x,s)\succ0
\quad\Longrightarrow\quad
T_n(x,s)+a^2e_ne_n^T\succ0.
\tag{25}
\]
In pivot language this is precisely
\[
\rho_1(x,s)>0,\dots,\rho_n(x,s)>0.
\tag{26}
\]
Thus the problem is exactly to remove the final positive rank-one boundary term \(a^2e_ne_n^T\), not pointwise but after sliding \(x\) inside the neighboring spectral gaps of \(A_n(a)\).

This can be summarized as the exact boundary-removal problem:
\[
\boxed{
\begin{minipage}{0.91\linewidth}
Assume \(T_n(x,s)+a^2e_ne_n^T\succ0\) at a point \(x\) lying between two consecutive eigenvalues of \(A_{n+1}(a)\).  Let \(\lambda_{n,k}(a)\) be the unique eigenvalue of \(A_n(a)\) lying in that interval.  Then one must prove that either
\[
T_n(y,s)\succ0
\quad\text{for some }y\in(\lambda_{n,k-1}(a),\lambda_{n,k}(a)),
\]
or
\[
T_n(y,s)\succ0
\quad\text{for some }y\in(\lambda_{n,k}(a),\lambda_{n,k+1}(a)).
\]
\end{minipage}}
\tag{27}
\]
This is equivalent to the localized barrier comparison
\[
\beta_{n+1,k}(a)
\le
\max\{\beta_{n,k-1}(a),\beta_{n,k}(a)\},
\tag{28}
\]
with the conventions \(\beta_{n,0}=\beta_{n,n}=0\).

At this point, I do not have a gap-free proof of the oscillation step (27).  The calculations above show that the missing argument is not a routine Riccati induction: the scalar recursion gives the exact finite-section criterion, but the desired implication requires a genuine global Sturm separation theorem for the \(2\times2\) block Jacobi pencil (12).  A correct proof must establish that boundary removal in (25) can always be compensated by moving \(x\) across the unique interlacing zero \(\lambda_{n,k}(a)\).  Without that separation theorem, the proof would be circular.

\end{document}